\documentclass[12pt,oneside]{book}

% ==========================================================
% METADATOS
% ==========================================================
\newcommand{\universidad}{Universidad de Buenos Aires}
\newcommand{\facultad}{Facultad de Derecho}
\newcommand{\materia}{Lectocomprensión Jurídica en Inglés}
\newcommand{\titulo}{Apuntes de Lectocomprensión Jurídica en Inglés}
\newcommand{\autor}{Isaac Edgar Camacho Ocampo}
\newcommand{\anio}{2026}

% ==========================================================
% PAQUETES
% ==========================================================
\usepackage[utf8]{inputenc}
\usepackage[T1]{fontenc}
% Babel no está disponible en esta instalación; se omite para evitar errores de compilación.

\usepackage[
    top=2.5cm,
    bottom=2.5cm,
    left=3cm,
    right=2.5cm,
    headheight=15pt
]{geometry}

\usepackage{amsmath}
\usepackage{amssymb}
\usepackage{mathtools}

\usepackage{graphicx}
\usepackage{float}
\usepackage{caption}
\usepackage{subcaption}

\usepackage{booktabs}
\usepackage{array}
\usepackage{longtable}
\usepackage{tabularx}

\usepackage{enumitem}

\usepackage{xcolor}
\usepackage[most]{tcolorbox}

\usepackage{fancyhdr}
\usepackage{ragged2e}

\usepackage[
    colorlinks=true,
    linkcolor=blue!50!black,
    citecolor=green!40!black,
    urlcolor=blue!60!black,
    pdftitle={\titulo},
    pdfauthor={\autor},
    pdfsubject={Apuntes universitarios},
    bookmarks=true
]{hyperref}

% ==========================================================
% CONFIGURACIÓN
% ==========================================================
\graphicspath{
    {./img/}
    {./graficos/}
    {./figuras/}
}

\setcounter{tocdepth}{2}

\setlength{\parindent}{0pt}
\setlength{\parskip}{0.5em}

\setlist[itemize]{
    topsep=0.3em,
    itemsep=0.2em
}

\setlist[enumerate]{
    topsep=0.3em,
    itemsep=0.2em
}

\clubpenalty=10000
\widowpenalty=10000

% ==========================================================
% COMANDOS
% ==========================================================
\newcommand{\abs}[1]{\left|#1\right|}
\newcommand{\norm}[1]{\left\lVert#1\right\rVert}
\newcommand{\vect}[1]{\mathbf{#1}}
\newcommand{\deriv}[2]{\frac{d#1}{d#2}}
\newcommand{\pderiv}[2]{\frac{\partial#1}{\partial#2}}

% Fila de pregunta/respuesta del cuadro Cornell
\newcommand{\cornellrow}[2]{\textbf{#1} & #2 \\ \addlinespace[0.4em]}

% ==========================================================
% ENTORNOS / CAJAS
% ==========================================================
\newtcolorbox{definition}[1]{
    enhanced,
    breakable,
    title={Definición: #1},
    fonttitle=\bfseries,
    colback=blue!3,
    colframe=blue!50!black,
    boxrule=0.8pt,
    arc=2mm
}

\newtcolorbox{important}{
    enhanced,
    breakable,
    title={Importante},
    fonttitle=\bfseries,
    colback=yellow!8,
    colframe=orange!70!black,
    boxrule=0.8pt,
    arc=2mm
}

\newtcolorbox{example}[1]{
    enhanced,
    breakable,
    title={Ejemplo: #1},
    fonttitle=\bfseries,
    colback=green!4,
    colframe=green!40!black,
    boxrule=0.8pt,
    arc=2mm
}

\newtcolorbox{formula}[1]{
    enhanced,
    breakable,
    title={Fórmula / Regla: #1},
    fonttitle=\bfseries,
    colback=purple!4,
    colframe=purple!50!black,
    boxrule=0.8pt,
    arc=2mm
}

\newtcolorbox{exercise}[1]{
    enhanced,
    breakable,
    title={Ejercicio: #1},
    fonttitle=\bfseries,
    colback=gray!5,
    colframe=gray!60!black,
    boxrule=0.8pt,
    arc=2mm
}

\newtcolorbox{note}{
    enhanced,
    breakable,
    title={Nota},
    fonttitle=\bfseries,
    colback=white,
    colframe=gray!50,
    boxrule=0.6pt,
    arc=2mm
}

% ==========================================================
% CONFIGURACIÓN FINAL
% ==========================================================
\pagestyle{fancy}
\fancyhf{}

\fancyhead[L]{\nouppercase{\leftmark}}
\fancyhead[R]{\materia}

\fancyfoot[C]{\thepage}

\renewcommand{\headrulewidth}{0.4pt}
\renewcommand{\footrulewidth}{0pt}

\fancypagestyle{plain}{
    \fancyhf{}
    \fancyfoot[C]{\thepage}
    \renewcommand{\headrulewidth}{0pt}
}

% ==========================================================
% DOCUMENTO
% ==========================================================
\begin{document}

% ----------------------------------------------------------
% PORTADA
% ----------------------------------------------------------
\begin{titlepage}

\thispagestyle{empty}

\begin{center}

\vspace*{1cm}

{\large \textsc{\universidad}}

\vspace{0.2cm}

{\large \textsc{\facultad}}

\vspace{3cm}

{\Huge \textbf{\titulo}}

\vspace{1cm}

{\LARGE Clase del 28 de octubre de 2026}

\vspace{0.8cm}

\textit{Comprender los textos con método es el primer paso para convertir el estudio en criterio profesional.}

\vspace{1.2cm}

\rule{0.70\textwidth}{0.5pt}

\vspace{0.5cm}

{\large Apuntes de estudio}

\vspace{0.5cm}

\rule{0.70\textwidth}{0.5pt}

\vfill

{\large \autor}

\vspace{0.3cm}

{\large \anio}

\end{center}

\vfill

\begin{flushleft}
\textit{Este es un modesto aporte a los estudiantes de ``Lectocomprensión Jurídica en Inglés'' no pretende sustituir el material oficial de la materia.}
\end{flushleft}

\end{titlepage}

% ----------------------------------------------------------
% ÍNDICE
% ----------------------------------------------------------
\tableofcontents

% ----------------------------------------------------------
% CONTENIDO
% ----------------------------------------------------------

% ==========================================================
\chapter[Utilidad profesional]{Utilidad profesional de la lectocomprensión jurídica}
% ==========================================================

Los contenidos de esta clase comprenden la corrección del parcial basado en el fallo \textit{Edmonson v. Leesville Concrete Co.} (Corte Suprema de los Estados Unidos, 1991), con los siguientes ejes: \textbf{acción estatal}, \textbf{recusaciones con y sin causa}, los \textbf{tres pasos de \textit{Batson}} (\textit{prima facie showing}, \textit{burden shifting}, \textit{judicial determination}), la \textbf{culpa concurrente} y el \textbf{\textit{certiorari}}. Además se trataron la modalidad de la segunda evaluación en tres fases y una reflexión sobre el uso responsable de la inteligencia artificial en el estudio y en las profesiones.

\section{¿Para qué sirve esto en la vida profesional?}

Los apuntes plantean que el trabajo con jurisprudencia en inglés es útil por tres razones:

\begin{enumerate}
    \item \textbf{Lectura metódica de fallos.} Leer fallos en inglés con método (partes, hechos, cuestión jurídica, decisión) es una habilidad directa para el derecho comparado, los contratos internacionales y la investigación de jurisprudencia extranjera.
    \item \textbf{Vocabulario técnico.} Términos como \textit{burden of proof}, \textit{remanded}, \textit{certiorari} y \textit{state action} aparecen en contratos, sentencias y doctrina que un abogado o traductor encuentra todos los días.
    \item \textbf{Uso de herramientas de IA con criterio propio.} Quien sabe hacer el trabajo manualmente es quien detecta los errores de la máquina.
\end{enumerate}

\begin{note}
Estructura de lectura de un fallo utilizada en la clase: \textbf{partes}, \textbf{tribunales}, \textbf{hechos}, \textbf{cuestión jurídica} y \textbf{decisión}.
\end{note}

% ==========================================================
\chapter[El fallo \textit{Edmonson}]{El fallo \textit{Edmonson v. Leesville Concrete Co.}}
% ==========================================================

\section{Precedentes y datos de citación}

El análisis parte de dos fallos de la Corte Suprema de los Estados Unidos:

\begin{itemize}
    \item \textbf{\textit{Batson v. Kentucky}}, 476 U.S. 79 (1986): precedente en materia penal que prohíbe usar recusaciones sin causa con motivación racial en la selección del jurado.
    \item \textbf{\textit{Edmonson v. Leesville Concrete Co.}}, 500 U.S. 614 (1991): la Corte Suprema aplica la doctrina de \textit{Batson} a los juicios civiles (decisión 6 a 3, voto del juez Kennedy).
\end{itemize}

\begin{note}
Los datos de citación fueron incorporados en el documento original de apuntes para identificar los fallos; en clase se los mencionó por su nombre.
\end{note}

\section{Partes, tribunales y hechos}

\subsection{Identificación de las partes}

En el análisis de un fallo no alcanza con mencionar los nombres de las partes: es necesario indicar cuál es la \textbf{actora} y cuál la \textbf{demandada}. En la fórmula «X \textit{versus} Y», el \textit{versus} indica quién es la actora.

\subsection{Tribunales intervinientes}

\begin{itemize}
    \item \textbf{Primera instancia:} \textit{US District Court for the Western District of Louisiana}. \textit{District} designa un tribunal de primera instancia; es federal porque lleva el \textit{US} adelante; y corresponde al distrito de Luisiana.
    \item \textbf{Última instancia:} la Corte Suprema federal de los Estados Unidos, a la que el caso llegó en apelación.
\end{itemize}

\subsection{Hechos principales}

Los hechos tienen dos componentes:

\begin{enumerate}
    \item \textbf{Hechos que originaron la causa (conflicto de fondo).} Un empleado afroamericano, trabajador de la construcción, sufrió lesiones en su lugar de trabajo (\textit{sustaining injuries at the job site}) por \textbf{negligencia}: no hubo dolo ni intención. Demandó por los daños.
    \item \textbf{Incidente en el proceso.} Al conformarse el jurado, la contraparte ejerció recusaciones sin causa sobre todos los posibles jurados afroamericanos, con el fin de que el jurado de primera instancia fuera lo más blanco posible. Esto permite sospechar el ejercicio de criterios discriminatorios.
\end{enumerate}

\section{Cuestión jurídica y acción estatal}

La \textbf{cuestión jurídica} consiste en determinar si es viable ejercer las recusaciones sin causa de esta manera y hasta dónde puede intervenir el Estado.

\begin{definition}{Acción estatal}
Cuando se habla de acción estatal se alude a \textbf{hasta dónde puede intervenir el Estado}. Los abogados de las partes tienen consagrado el derecho a ejercer recusaciones, que son lícitas y legítimas. El problema surge cuando, \textit{prima facie}, se observa discriminación: allí deben intervenir el Estado y las autoridades judiciales, porque se trata de un proceso supervisado por la justicia y por el juez.
\end{definition}

Se explicó que las recusaciones sin causa de partes privadas constituyen acción estatal porque lo que está en juego son \textbf{facultades otorgadas por ley}, ejercidas \textbf{dentro de una estructura pública} (un juzgado, un procedimiento) y \textbf{bajo la supervisión de un juez}.

\subsection{Cláusula de igual protección}

La cláusula de igual protección de la Constitución fue mencionada como un dato adicional de la respuesta (su inclusión no era indispensable).

\begin{note}
Texto citado en los apuntes: «\textit{\ldots nor deny to any person within its jurisdiction the equal protection of the laws.}» (Enmienda XIV, sección 1, Constitución de los Estados Unidos). Según la nota del documento, por tratarse de un tribunal federal, en \textit{Edmonson} la garantía se aplica a través del componente de igual protección de la Quinta Enmienda. La docente se refiere a ella como «cláusula de igual protección».
\end{note}

\section{Recusaciones con causa y sin causa}

Los \textit{challenges} pueden ser con causa o sin causa. Para que la respuesta sea completa debe indicarse que unas son \textbf{ilimitadas} y otras \textbf{limitadas}.

\begin{table}[H]
\centering
\renewcommand{\arraystretch}{1.3}
\begin{tabularx}{\textwidth}{
    >{\raggedright\arraybackslash}p{0.18\textwidth}
    >{\raggedright\arraybackslash}X
    >{\raggedright\arraybackslash}X
}
\toprule
 & \textbf{Recusación con causa} (\textit{challenge for cause}) & \textbf{Recusación sin causa} (\textit{peremptory challenge}) \\
\midrule
\textbf{Fundamentación} & Se justifica y se fundamenta. & No hace falta decir por qué: es la esencia de esta recusación. \\
\textbf{Cantidad} & Ilimitadas: todas las que haga falta. & Limitadas: «tres o cinco, depende del estado». \\
\textbf{Control} & Es fundada, por lo tanto controlable. & Puede ser cuestionada si hay indicios de discriminación; entonces se invierte la carga y quien recusó debe dar otras razones. \\
\bottomrule
\end{tabularx}
\end{table}

\section{Los tres pasos de \textit{Batson}}

La Corte estableció tres pasos para determinar si existió discriminación racial en la selección del jurado según \textit{Batson v. Kentucky}.

\begin{table}[H]
\centering
\renewcommand{\arraystretch}{1.3}
\begin{tabularx}{\textwidth}{
    >{\raggedright\arraybackslash}p{0.30\textwidth}
    >{\raggedright\arraybackslash}X
}
\toprule
\textbf{Paso (en inglés)} & \textbf{Explicación} \\
\midrule
\textbf{1. \textit{Prima facie showing} of discrimination} & Una demostración que se pueda comprobar. \textit{Prima facie} no es inglés, es latín: «a primera vista», algo manifiesto, evidente; que sin un análisis demasiado profundo se pueda demostrar o presuponer que hubo discriminación racial. \\
\addlinespace
\textbf{2. \textit{Burden shifting} to the striking party} & \textit{Shift} es dar vuelta, cambiar; \textit{burden} es carga (de la prueba o de justificación). Es dar vuelta la carga de la prueba, de modo que no recaiga en la parte a la que normalmente le hubiera correspondido, sino en la otra. \\
\addlinespace
\textbf{3. \textit{Judicial determination} of the justification's validity} & Una resolución judicial que determine la validez del planteo. Es el juez quien evalúa las pruebas presentadas y determina si quien objeta efectivamente pudo probar que hubo discriminación racial. No sucede de manera espontánea. \\
\bottomrule
\end{tabularx}
\end{table}

\begin{note}
Sobre \textit{burden shifting}: se recordó lo visto con las cláusulas nominales, donde la lectura se inicia «de atrás para adelante». Aquí ocurre lo mismo.
\end{note}

\begin{definition}{\textit{Prima facie}}
Expresión latina que significa «a primera vista». En el proceso no hace falta mostrar los hechos sin lugar a dudas: una primera evaluación puede demostrar que existen indicios para probar una discriminación racial.
\end{definition}

\begin{definition}{\textit{Burden shifting}}
Inversión de la carga de la prueba: quien alega no haber cometido la discriminación debe demostrarlo, dando razones distintas de la raza para sus recusaciones.
\end{definition}

\begin{important}
Esquema del razonamiento: (1) recusaciones sin causa que siguen un patrón $\rightarrow$ (2) indicio \textit{prima facie} de discriminación $\rightarrow$ (3) se invierte la carga de la prueba (\textit{burden shifting}) $\rightarrow$ (4) quien recusó debe dar razones neutrales, no raciales $\rightarrow$ (5) el juez evalúa si son válidas (\textit{judicial determination}).
\end{important}

\section{Resolución de la Corte y \textit{remanded}}

La Corte resolvió, en un fallo dividido de seis a tres (no unánime), que las recusaciones sin causa por parte de particulares \textbf{sí son acción estatal} y quedan sujetas a la prohibición de discriminación racial, aplicando la doctrina del caso \textit{Batson v. Kentucky}.

Hacia el final del texto del fallo, antes de la referencia a \textit{instructions}, aparece el término \textit{remanded}.

% TODO: contenido incompleto o ambiguo en las notas originales (fragmento poco claro en la grabación tras «to determine\ldots»)

\begin{definition}{\textit{Remanded}}
Cuando una causa llega en apelación, el tribunal puede \textbf{confirmar} la sentencia de instancia anterior, \textbf{revocarla} o \textbf{devolver} las actuaciones. \textit{Remanded} corresponde a esta tercera posibilidad.
\end{definition}

Explicación textual de la docente sobre el alcance de la devolución: «Básicamente dijo: ``Tomá y volvé a fijarte si no lo pudieron demostrar''. Que den otras razones, que no estén basadas en raza, para estas recusaciones. Bueno, okay, que lo demuestre. Por eso lo de \textit{shifting}: la inversión de la carga. Tiene que demostrar que no fue así.»

\section{Historia procesal, \textit{certiorari} y valor del precedente}

La historia procesal no se pedía en ninguna consigna, pero sirve para comprender el caso. Un tribunal del Quinto Circuito dijo inicialmente que el precedente de \textit{Batson} se aplicaba y reenvió la causa para que se presentaran justificaciones neutrales (¿por qué se recusó a esos jurados?, no por raza). Luego, ese mismo tribunal, en pleno, decidió revocar: sostuvo que el principio de \textit{Batson} no se aplicaba a las causas civiles, sino a lo penal. Entonces intervino la Corte Suprema, que otorgó el \textit{certiorari}.

La Corte tuvo que resolver si la cláusula constitucional de protección igualitaria y el principio de no discriminación en la elección de jurados también se aplicaban a las causas civiles.

\begin{definition}{\textit{Certiorari}}
Recurso mediante el cual la Corte Suprema elige discrecionalmente qué casos revisar. En clase se lo comparó con el recurso de avocación y con el recurso extraordinario federal, y en algunos aspectos con el \textit{per saltum} (cuando se salta una instancia y se va directamente a la Corte Suprema por ciertos motivos). Es un recurso que no se otorga en cualquier circunstancia.
\end{definition}

\begin{table}[H]
\centering
\renewcommand{\arraystretch}{1.3}
\begin{tabularx}{\textwidth}{
    >{\raggedright\arraybackslash}p{0.19\textwidth}
    >{\raggedright\arraybackslash}p{0.25\textwidth}
    >{\raggedright\arraybackslash}X
}
\toprule
\textbf{Instancia} & \textbf{Tribunal} & \textbf{Decisión} \\
\midrule
\textbf{Primera instancia} & \textit{US District Court, Western District of Louisiana} & Jurado; recusaciones sin causa; se determina 80\,\% de culpa concurrente del actor. \\
\addlinespace
\textbf{Apelación} & Tribunal del Quinto Circuito & Primero aplica \textit{Batson} y devuelve la causa; luego, en pleno, revoca: \textit{Batson} no rige en lo civil. \\
\addlinespace
\textbf{Última instancia} & Corte Suprema de los Estados Unidos & \textit{Certiorari}. Fallo 6 a 3: sí hay acción estatal; \textit{Batson} se aplica a lo civil; \textit{remanded}. \\
\bottomrule
\end{tabularx}
\end{table}

\textbf{Valor del precedente.} Lo importante era sentar jurisprudencia y dejar establecido el criterio, lo cual se vincula con la dinámica del proceso anglosajón: los precedentes se superponen como «capas geológicas». En este caso hay dos en juego: \textit{Batson} y \textit{Edmonson v. Leesville}, que deja un nuevo precedente, al establecer que sí constituye acción estatal y que el juez debe intervenir, también en casos civiles, para analizar la cuestión de la discriminación.

\section{Cuadro Cornell: el fallo \textit{Edmonson}}

\begin{longtable}{@{}>{\raggedright\arraybackslash}p{0.27\textwidth}>{\raggedright\arraybackslash}p{0.69\textwidth}@{}}
\toprule
\textbf{Preguntas / palabras clave} & \textbf{Notas / respuestas} \\
\midrule
\endfirsthead
\toprule
\textbf{Preguntas / palabras clave} & \textbf{Notas / respuestas} \\
\midrule
\endhead
\cornellrow{¿Qué tribunales intervinieron en el caso?}{\textit{US District Court for the Western District of Louisiana} (primera instancia federal, distrito de Luisiana); Tribunal del Quinto Circuito (apelación); Corte Suprema de los Estados Unidos (última instancia).}
\cornellrow{¿Cuál fue el conflicto de fondo y cuál la cuestión jurídica?}{De fondo: lesiones sufridas por un empleado afroamericano en su lugar de trabajo por negligencia (sin dolo). Jurídica: si es viable ejercer recusaciones sin causa de esta manera y hasta dónde puede intervenir el Estado.}
\cornellrow{¿Qué es la acción estatal y por qué las recusaciones sin causa de particulares la constituyen?}{Es la intervención del Estado. Las recusaciones sin causa de particulares son acción estatal porque ocurren dentro de una estructura pública (un juzgado), con facultades otorgadas por ley y bajo la supervisión de un juez.}
\cornellrow{¿En qué difieren las recusaciones con causa y sin causa?}{Con causa: ilimitadas, fundadas y controlables. Sin causa (\textit{peremptory}): limitadas (tres o cinco, según el estado) y sin necesidad de explicar, salvo que haya indicio de discriminación; en ese caso se invierte la carga y quien recusó debe dar otras razones.}
\cornellrow{¿Cuáles son los tres pasos de \textit{Batson}?}{1) \textit{Prima facie showing} of discrimination; 2) \textit{burden shifting} to the striking party; 3) \textit{judicial determination} of the justification's validity.}
\cornellrow{¿Qué significan \textit{prima facie} y \textit{burden shifting}?}{\textit{Prima facie}: «a primera vista»; indicios suficientes, sin análisis profundo, para presumir discriminación racial. \textit{Burden shifting}: inversión de la carga de la prueba; quien recusó debe dar razones neutrales, no raciales.}
\cornellrow{¿Qué resolvió la Corte Suprema y qué significa \textit{remanded}?}{En un fallo 6 a 3, que las recusaciones sin causa de particulares son acción estatal sujeta a la prohibición de discriminación racial, aplicando \textit{Batson}. \textit{Remanded} es devolver la causa al tribunal inferior, tercera posibilidad del tribunal de apelación junto con confirmar o revocar.}
\cornellrow{¿Qué es el \textit{certiorari} y por qué importa el precedente?}{Recurso discrecional ante la Corte Suprema, comparable al extraordinario federal. Importa porque el precedente se acumula como capas geológicas: \textit{Batson} (penal) y \textit{Edmonson} (civil).}
\midrule
\multicolumn{2}{@{}>{\raggedright\arraybackslash}p{0.96\textwidth}@{}}{\textbf{Resumen:} \textit{Edmonson v. Leesville Concrete Co.} extendió a los juicios civiles la doctrina de \textit{Batson}. Las recusaciones sin causa de particulares son acción estatal; ante un indicio \textit{prima facie} de discriminación se invierte la carga de la prueba, quien recusó debe dar razones neutrales y un juez determina su validez. La Corte, por 6 a 3, otorgó el \textit{certiorari} y devolvió la causa (\textit{remanded}).} \\
\bottomrule
\end{longtable}

% ==========================================================
\chapter{Corrección del parcial}
% ==========================================================

El parcial se basó en el fallo \textit{Edmonson v. Leesville Concrete Co.} Las consignas fueron corregidas con los criterios que se detallan a continuación.

\section{Verdadero o falso y justificación}

En esta consigna debían justificarse \textbf{solo las falsas}. La consigna está completa únicamente si se justifican las falsas: si falta la justificación, la respuesta no alcanza. Las falsas eran la primera y la última.

\subsection{Ítem A}

Parte de la justificación consistía en que no se trató de un fallo unánime, sino dividido, seis a tres, con el juez Kennedy y tres votos en disidencia.

\subsection{Ítem D}

\begin{exercise}{Ítem D (verdadero o falso)}
\textbf{Enunciado:} «La Corte estableció que las recusaciones sin causa realizadas por partes privadas no constituyen acción estatal porque no intervienen funcionarios públicos.»

\textbf{Corrección:} era \textbf{falsa}. La Corte dijo que sí constituye acción estatal, porque están en juego competencias o facultades otorgadas por ley, dentro de una estructura pública (un juzgado, un procedimiento), bajo la supervisión de un juez. Por lo tanto, hay acción estatal e intervención del Estado.
\end{exercise}

\section{Opción múltiple y análisis de distractores}

La única opción correcta era la \textbf{B}.

\begin{exercise}{Opción B (correcta)}
«El uso de recusaciones sin causa por particulares en juicios civiles se realiza bajo prestación legal dentro del tribunal y bajo supervisión judicial, constituyendo una acción estatal.»
\end{exercise}

Las demás opciones contenían algún aspecto que impedía que fueran correctas:

\begin{table}[H]
\centering
\renewcommand{\arraystretch}{1.3}
\begin{tabularx}{\textwidth}{
    >{\raggedright\arraybackslash}p{0.12\textwidth}
    >{\raggedright\arraybackslash}X
}
\toprule
\textbf{Opción} & \textbf{Qué decía y por qué no era correcta} \\
\midrule
\textbf{A} & Decía que se trataba estrictamente de un caso de discriminación laboral del demandante, ciudadano afroamericano. No había habido discriminación laboral: el conflicto laboral era otro. El tema de la discriminación se presentó en la elección del jurado. \\
\addlinespace
\textbf{C} & Decía que el tribunal de primera instancia incurrió en error grave al dictar instrucciones incorrectas al jurado. Muchos se apoyaron en esto. El problema no estuvo en las instrucciones, sino en las recusaciones: en ellas se estaba filtrando la cuestión de la discriminación, que no debería suceder. \\
\addlinespace
\textbf{D} & Decía que el precedente \textit{Batson} había sido ampliado expresamente por el Congreso mediante una ley aprobada. No tenía relación con el caso. \\
\bottomrule
\end{tabularx}
\end{table}

% TODO: contenido incompleto o ambiguo en las notas originales (la docente no pudo precisar de dónde surgió la confusión de varios estudiantes con las instrucciones al jurado)

\begin{note}
La docente indicó que no sabía con precisión qué parte del texto generó la confusión acerca de las instrucciones al jurado.
\end{note}

\section{Los tres pasos de \textit{Batson}}

La consigna pedía, en forma breve, los tres pasos establecidos por la Corte para determinar si existió discriminación racial en la selección del jurado según \textit{Batson v. Kentucky}. Los tres pasos estaban mencionados en el texto, en tres puntos: \textit{prima facie showing}, \textit{burden shifting} y \textit{judicial determination} (véase el capítulo anterior).

\section{Cuadro de partes, tribunales y hechos}

\begin{itemize}
    \item \textbf{Partes:} no alcanza con mencionar los nombres de las empresas o de las partes; debe indicarse expresamente cuál es la actora y cuál la demandada. El \textit{versus} marca quién es la actora, pero nombrar a las partes sin indicar su rol no es suficiente.
    \item \textbf{Tribunales:} último, la Corte Suprema de los Estados Unidos; primero, \textit{US District Court for the Western District of Louisiana}.
    \item \textbf{Hechos:} se distinguen los hechos de la causa (lesiones por negligencia de un empleado afroamericano) y las recusaciones sin causa que permiten sospechar criterios discriminatorios.
\end{itemize}

\section{Consigna sobre los \textit{challenges}}

La consigna sobre los \textit{challenges} requería distinguir recusación con causa y sin causa. Para que la respuesta fuera completa convenía incluir que las primeras son ilimitadas y las segundas limitadas: las que son con causa pueden ser todas las que hagan falta, porque se justifican; las que son sin causa son limitadas, y pueden ser cuestionadas por la otra parte, de modo que luego deba analizarse si hubo una discriminación oculta.

\section{Cuadro Cornell: corrección del parcial}

\begin{longtable}{@{}>{\raggedright\arraybackslash}p{0.27\textwidth}>{\raggedright\arraybackslash}p{0.69\textwidth}@{}}
\toprule
\textbf{Preguntas / palabras clave} & \textbf{Notas / respuestas} \\
\midrule
\endfirsthead
\toprule
\textbf{Preguntas / palabras clave} & \textbf{Notas / respuestas} \\
\midrule
\endhead
\cornellrow{¿Qué exigía la consigna de verdadero o falso?}{Justificar solo las falsas. Sin justificación, la respuesta está incompleta.}
\cornellrow{¿Por qué \textit{Edmonson} no fue un fallo unánime?}{Fue un fallo dividido, seis a tres: el juez Kennedy y tres votos en disidencia.}
\cornellrow{¿Por qué era falso el ítem D?}{Porque la Corte dijo que las recusaciones sin causa de particulares sí constituyen acción estatal: se trata de facultades otorgadas por ley, dentro de un juzgado y bajo supervisión judicial.}
\cornellrow{¿Por qué las opciones A, C y D de la opción múltiple eran incorrectas?}{A: no hubo discriminación laboral; la discriminación se presentó en la elección del jurado. C: el problema no estuvo en las instrucciones sino en las recusaciones. D: no hubo ampliación de \textit{Batson} por ley del Congreso.}
\cornellrow{¿Qué debía indicarse sobre las partes en el cuadro?}{Cuál es la actora y cuál la demandada; mencionar solo los nombres no alcanza.}
\midrule
\multicolumn{2}{@{}>{\raggedright\arraybackslash}p{0.96\textwidth}@{}}{\textbf{Resumen:} la corrección enfatizó la completitud de las respuestas: justificar las falsas, identificar el rol de las partes y distinguir con precisión las recusaciones con y sin causa. La opción correcta fue la B (acción estatal por actuar bajo prestación legal y supervisión judicial).} \\
\bottomrule
\end{longtable}

% ==========================================================
\chapter[Debates sobre el fallo]{Debates a partir del fallo}
% ==========================================================

\section{Imparcialidad del juez e instrucciones al jurado}

Un estudiante advirtió de dónde surgía la mención de las instrucciones al jurado en la conclusión del texto: las cortes deben guiar a los jurados para que deliberen de modo justo e imparcial y valoren la prueba, siempre que los jurados hayan sido elegidos de manera adecuada y legítima.

La docente señaló que el mismo juez que permite recusar a un jurado por ser afroamericano es quien después instruye al jurado. Esto pone en duda, en última instancia, la \textbf{integridad de todo el proceso}. En una apelación, la parte que pierda podría sostener: «Si el juez vio esto y no hizo nada, ¿cómo estoy seguro o segura de que no tiene un sesgo racial el propio juez al momento de las instrucciones?». El texto lo menciona hacia el final con la expresión \textit{procedural integrity}, que reúne toda la integridad procesal.

Se planteó además otra circunstancia: durante la producción de prueba, con el jurado ya elegido, el juez admite o rechaza objeciones sobre medios de prueba y dirige el proceso, por lo que existen numerosas instancias en las que puede cuestionarse su imparcialidad.

\begin{note}
Lo que el fallo pone en tela de juicio sobre el final: si durante las recusaciones y la conformación del jurado el juez hizo la vista gorda, ¿qué garantía hay cuando llegue el momento de las instrucciones? En palabras de la docente: «Plantea la duda; tira ahí la piedrita.»
\end{note}

\section{El conflicto de fondo}

Ante la consulta de un estudiante, se aclaró que el conflicto de fondo es una \textbf{negligencia laboral}: el trabajador sufrió lesiones en su lugar de trabajo. Cuando se pregunta por los hechos principales de la causa, se alude a ese daño. La demanda se origina por eso; es negligente o culposo, sin intención. En el curso del proceso, al conformarse el jurado, las recusaciones sin causa excluyeron a todas las personas de color de los posibles jurados.

\section{El indicio y la inversión de la carga de la prueba}

Cuando la recusación es sin causa, es difícil demostrar que existe un patrón o una motivación, porque no hay que explicar por qué se recusa.

\begin{example}{Patrón de recusaciones}
Si se recusa al jurado número cuatro, luego al siete y luego al nueve, sin necesidad de dar razones, y esos tres jurados son precisamente las tres personas negras, y una de las partes del juicio tiene esa característica o pertenece a esa comunidad, existe un \textbf{indicio}.
\end{example}

\textit{Prima facie} significa que no hace falta mostrar sin lugar a dudas la discriminación: una primera evaluación puede evidenciar indicios. Eso avala en cierta medida el reclamo y es entonces cuando se \textbf{invierte la carga de la prueba}: debe demostrar quien alega no haber discriminado. No se exige a quien reclama probar que el otro discriminó, sino que el recusante diga, por ejemplo, «lo recusé porque\ldots» y aporte otro motivo.

Esto rompe, en parte, la dinámica de la recusación sin causa, cuya esencia es no tener que fundamentar. Pero la obligación de explicar surge \textbf{únicamente por el patrón}, es decir, porque siempre se recusaba a las mismas personas y hubo un indicio o sospecha de discriminación. En ausencia de ello no hay que decir nada: se recusa y punto (tres o cinco recusaciones, según el estado).

\section{El precedente y el interés de las partes}

Un estudiante observó que la recusación ni siquiera era necesaria, pues el jurado que había quedado ya había fallado a favor del empleado; por lo tanto, se trataba de marcar un precedente. La docente coincidió: es marcar el precedente y abrir un tema espinoso para el sistema anglosajón y para la selección de jurados, donde existe una laguna latente.

\section{Monto de la indemnización y gravedad del daño}

En la discusión se planteó lo siguiente:

\begin{itemize}
    \item Un estudiante mencionó la cifra de 18.000 dólares en el contexto de un sistema de salud privado y sugirió buscar su equivalente a valores actuales, ya que el fallo es de 1991. La docente comentó que había pensado que la cifra era actual.
    \item El texto solo indica que el empleado de la construcción sufrió lesiones en el lugar de trabajo (\textit{sustaining injuries at the job site}); no precisa el accidente ni sus efectos. Un estudiante conjeturó una maquinaria mal operada o con mal mantenimiento por parte de la empresa; se reconoció que no se sabe más.
    \item La docente señaló que, si la lesión era incapacitante o dejaba consecuencias a futuro para conseguir otro trabajo, el monto podía resultar muy bajo. Un estudiante conjeturó que por eso se habría pedido la renovación del jurado.
\end{itemize}

La docente subrayó que lo importante era la existencia de un precedente, \textit{Batson}, donde se sienta el principio sobre cómo evaluar ciertas pautas. Añadió que, aunque se hubiera fallado a favor del empleado, el planteo podía deberse a que el monto resultó bajo. Observó también que \textit{Batson} era un fallo penal y que no podía mezclarse con lo civil, pues no rigen los mismos principios para la selección o recusación del jurado.

% TODO: contenido incompleto o ambiguo en las notas originales (la docente indica que lo penal y lo civil no se mezclan; un estudiante concluye que el fallo muestra que pueden compararse y a veces usar las mismas bases)

Un estudiante concluyó que el fallo muestra que lo civil y lo penal pueden compararse y que a veces pueden utilizarse las mismas bases.

\section{La culpa concurrente}

Un estudiante encontró en el texto el motivo de la reducción de la indemnización: el jurado determinó que Edmonson tuvo un \textbf{80\,\% de responsabilidad}. La docente lo identificó como culpa concurrente y observó que ochenta es un porcentaje muy alto. El estudiante interpretó que el empleado se había equivocado en la maniobra de una máquina.

\begin{definition}{Culpa concurrente (\textit{concurrent fault})}
La indemnización se redujo porque el jurado atribuyó a Edmonson el 80\,\% de la responsabilidad en el accidente.
\end{definition}

\section{Cuadro Cornell: debates sobre el fallo}

\begin{longtable}{@{}>{\raggedright\arraybackslash}p{0.27\textwidth}>{\raggedright\arraybackslash}p{0.69\textwidth}@{}}
\toprule
\textbf{Preguntas / palabras clave} & \textbf{Notas / respuestas} \\
\midrule
\endfirsthead
\toprule
\textbf{Preguntas / palabras clave} & \textbf{Notas / respuestas} \\
\midrule
\endhead
\cornellrow{¿Por qué se discutió la imparcialidad del juez?}{Si el juez no interviene ante recusaciones discriminatorias, podría dudarse de su imparcialidad al dar instrucciones al jurado y durante el proceso: se afecta la integridad procesal (\textit{procedural integrity}).}
\cornellrow{¿Cuándo se invierte la carga de la prueba?}{Cuando hay un indicio \textit{prima facie} de discriminación, derivado del patrón de recusaciones sin causa. Entonces quien recusó debe dar otras razones; sin patrón, no tiene que explicar nada.}
\cornellrow{¿Qué sentido tuvo marcar un precedente según la discusión?}{Abrir un tema espinoso para el sistema anglosajón y la selección de jurados, aunque el jurado ya hubiera fallado a favor del empleado.}
\cornellrow{¿Por qué se redujo la indemnización de Edmonson?}{Porque el jurado le atribuyó el 80\,\% de la responsabilidad: culpa concurrente.}
\midrule
\multicolumn{2}{@{}>{\raggedright\arraybackslash}p{0.96\textwidth}@{}}{\textbf{Resumen:} el debate mostró que las recusaciones discriminatorias afectan la integridad del proceso y la confianza en la imparcialidad del juez; que la carga de la prueba se invierte solo ante un patrón que genere indicio; que el caso importa como precedente; y que la reducción del monto se explica por la culpa concurrente (80\,\%).} \\
\bottomrule
\end{longtable}

% ==========================================================
\chapter[Segunda evaluación]{Segunda evaluación: cronograma y modalidad}
% ==========================================================

\section{Cronograma}

\begin{table}[H]
\centering
\renewcommand{\arraystretch}{1.3}
\begin{tabularx}{\textwidth}{
    >{\raggedright\arraybackslash}p{0.33\textwidth}
    >{\raggedright\arraybackslash}X
}
\toprule
\textbf{Fecha} & \textbf{Actividad} \\
\midrule
Miércoles 28 de octubre & Clase de hoy: devolución y corrección del primer parcial. \\
Martes siguiente (3 de noviembre) & Comienza la unidad de Familia (aproximadamente dos textos). \\
Miércoles 4 de noviembre & Recuperatorio del primer parcial. En principio sin clase; a confirmar por Classroom. \\
Miércoles 11 de noviembre & Segunda evaluación, fase 1: presencial y grupal. \\
Martes 17 de noviembre & Fase 2: consigna breve domiciliaria (2 o 3 días). \\
Miércoles 18 de noviembre & Fase 3: presencial y grupal. Cierre de la evaluación y del cuatrimestre. \\
Miércoles 25 de noviembre & Fecha reservada: recuperación oral individual para ausentes. \\
\bottomrule
\end{tabularx}
\end{table}

Según lo explicado, el recuperatorio del miércoles 4 se organizaría sin clase en principio, porque resulta incómodo dictar clase mientras hay estudiantes rindiendo; la decisión no estaba definida y se comunicaría por Classroom.

\section{Modalidad en tres fases}

La segunda evaluación sigue siendo un trabajo \textbf{grupal}, pero \textbf{presencial}, dividido en tres fases sobre un mismo texto. El examen es uno solo, dividido en tres partes, y en las tres participa el mismo grupo:

\begin{enumerate}
    \item \textbf{Fase 1 (presencial).} Se recibe el texto y se responde en grupo a determinadas consignas. Un grupo de cuatro entrega un solo trabajo, que la docente se lleva.
    \item \textbf{Fase 2 (domiciliaria).} Con el mismo texto y la misma actividad, se completa una consigna muy corta, con dos o tres días para resolverla. Se entrega un solo trabajo por grupo, enviado por mail por uno de los integrantes, con el nombre de todos. En principio no requiere asistir a clase.
    \item \textbf{Fase 3 (presencial).} Se vuelve a trabajar en el aula, con los mismos integrantes del grupo, con el mismo texto y las mismas consignas que la primera. Con ella se cierra la evaluación y el cuatrimestre.
\end{enumerate}

\subsection{Conformación de los grupos}

Los grupos se arman el mismo día de la evaluación, con los presentes y según como se sienten los estudiantes, para evitar grupos incompletos si alguien falta. Por eso se recomienda llegar temprano.

\subsection{Ausencias y recuperación}

\begin{important}
Quien esté ausente en alguna de las etapas presenciales de la segunda evaluación no tiene más instancias de recuperación que la reservada para fin de noviembre (calculada para el miércoles 25): un examen \textbf{oral e individual}, solo con la docente. La nota es grupal para quienes rinden en grupo.
\end{important}

\subsection{Material permitido}

La modalidad es igual que la del parcial: se puede traer todo el material, pero \textbf{sin IA y sin dispositivos}. Se aconseja traer apuntes y diccionario. El material de consulta corresponde a los temas ya trabajados. En la parte domiciliaria no se monitorea cómo se organizan los grupos. En el recuperatorio podría haber alguna variación en las consignas, con consulta de algunos de los temas trabajados (en el parcial se utilizó la modalidad de verdadero o falso).

\section{Fundamentos del cambio y contenido}

La evaluación seguirá siendo grupal, pero se modificaron las consignas y los temas de la segunda parte: se verán textos de familia y la evaluación versará principalmente sobre ellos. Los conocimientos empiezan a estar escalonados, por lo que no queda nada afuera: el texto de \textit{Edmonson} era contractual, pero también procesal y constitucional. «No son áreas estancas.»

\textbf{Razones del cambio de modalidad}, expresadas por la docente:

\begin{itemize}
    \item Con la modalidad anterior, un grupo de cuatro personas tenía una semana para hacer un trabajo en el que participaban, con suerte, dos integrantes que se ponían el trabajo al hombro, y los otros dos no. Resultaba una instancia que no evaluaba nada.
    \item La propuesta de elegir el tema permitía investigar a quien se lo tomaba en serio, pero quien no, exponía todo de memoria, y eso no servía.
    \item Antes era muy difícil individualizar quién trabajaba, quién no y quién necesitaba ayuda. Hacerlo en el aula permite intervenir, preguntar y orientar, para que trabajen todos y no se acumule trabajo en casa.
    \item También se cambió, según la docente, para frenar el avance de la inteligencia artificial, que está descontrolada.
\end{itemize}

\section{Cuadro Cornell: segunda evaluación}

\begin{longtable}{@{}>{\raggedright\arraybackslash}p{0.27\textwidth}>{\raggedright\arraybackslash}p{0.69\textwidth}@{}}
\toprule
\textbf{Preguntas / palabras clave} & \textbf{Notas / respuestas} \\
\midrule
\endfirsthead
\toprule
\textbf{Preguntas / palabras clave} & \textbf{Notas / respuestas} \\
\midrule
\endhead
\cornellrow{¿Cómo será la segunda evaluación?}{Grupal, en tres fases con el mismo texto: presencial (11/11), domiciliaria breve (17/11) y presencial (18/11). El contenido principal es Familia.}
\cornellrow{¿Qué ocurre con quien está ausente en una fase presencial?}{Debe rendir un examen oral e individual con la docente (fecha reservada: miércoles 25 de noviembre); no hay más instancias de recuperación.}
\cornellrow{¿Qué material se puede usar?}{Todo el material y diccionario, pero sin IA y sin dispositivos en las instancias presenciales.}
\cornellrow{¿Por qué cambió la modalidad?}{Porque el trabajo grupal anterior no permitía saber quién trabajaba; hacerlo en el aula permite acompañar a todos, evitar la acumulación de trabajo en casa y frenar el avance de la IA.}
\midrule
\multicolumn{2}{@{}>{\raggedright\arraybackslash}p{0.96\textwidth}@{}}{\textbf{Resumen:} la segunda evaluación es un único examen grupal dividido en tres fases (dos presenciales y una domiciliaria breve) sobre un mismo texto, centrado en Familia. Los ausentes rinden oral e individualmente el 25 de noviembre. El cambio busca que todos trabajen y permitir el acompañamiento docente.} \\
\bottomrule
\end{longtable}

% ==========================================================
\chapter[IA y juicio profesional]{Inteligencia artificial y juicio profesional}
% ==========================================================

\section{Uso responsable de la inteligencia artificial en el estudio}

La docente planteó que la inteligencia artificial, así como ayuda, también anula aspectos: se dejan de realizar operaciones que antes se hacían, porque uno se apoya en la herramienta y deja de pensar. Trasladar al ámbito educativo el uso cotidiano (recordatorios, listados, organización) y delegarle todo implica un riesgo: profesionales que pasaron por la facultad resolviendo todo con IA, como un médico o un arquitecto.

\begin{example}{Anécdota sobre la información que maneja una herramienta}
Una colega probó el prompt «escribime, en argentino, lo que sabés de mí». La respuesta incluyó en qué trabaja y para qué usa la herramienta, además de rasgos de personalidad y datos de sus hijos. Otra herramienta que usa menos no supo qué contestar. La docente concluyó que había una con la que estaba hablando demasiado.
\end{example}

La postura expresada no es contraria al uso: la docente la utiliza mucho para preparar clases, pero lo que generó tenía falencias (aspectos de la docencia, de la universidad y del perfil de alumno que no podía contemplar), por lo que fue necesario revisarlo y modificarlo. Lo central es \textbf{no perder la capacidad de cuestionar} el resultado ni la de \textbf{tomar la última decisión}.

\begin{important}
La IA es una herramienta: puede usarse, pero no debe perderse la capacidad de cuestionar, de verificar el resultado y de tomar la decisión final.
\end{important}

\subsection{La perspectiva de los estudiantes}

Los estudiantes plantearon que es un desafío ser más inteligente que la IA, y que la necesidad de ahorrar tiempo (por la cantidad de materias y trabajos) lleva a usarla. Se preguntaron si es la IA o la falta de tiempo lo que empobrece el pensamiento, ya que una herramienta que resuelve en poco tiempo lo que antes insumía mucho resulta atractiva. La docente respondió que, como se había dicho, es una herramienta.

\section{Las profesiones frente a la IA}

\begin{example}{Traducción}
La docente señaló que la traducción automática afectó fuertemente al traductorado, porque pocas empresas y clientes conocen el riesgo de dejar la traducción en manos de la IA. Ella la usa porque agiliza los tiempos, pero sostiene que no podría corregir ni detectar sus errores si no hubiera aprendido a traducir ni estudiado dos sistemas jurídicos comparados. Al traducir un balance extenso, advierte errores de sentido: la herramienta omite párrafos o partes. Corregir no es solo revisar: es mantener todas las antenas paradas para detectar errores.
\end{example}

\begin{itemize}
    \item \textbf{Diseño:} un estudiante comentó que los diseñadores deben de estar afectados, pues todos los \textit{flyers} resultan iguales. La docente mencionó la creación de logos, que antes era una actividad creativa con diálogo con el cliente sobre colores y qué proyectar de la marca. Un estudiante señaló que ahora se hace con el chat, y se paga. La docente respondió que la capacidad profesional sigue siendo otra cosa.
    \item \textbf{Derecho:} se mencionaron fallos falsos y redacciones todas iguales, con errores que a veces pasan desapercibidos.
\end{itemize}

\subsection{El aspecto humano del derecho}

La docente señaló que en cuestiones de familia, o en el caso de quien tuvo el accidente y tenía el 80\,\% de la responsabilidad (culpa concurrente), el aspecto humano y sus matices no son percibidos por la máquina. Sostuvo que la IA debe estar en el aula, porque está en la vida cotidiana, pero que no puede perderse la capacidad de aprender a hacer las cosas de manera analógica. Planteó el interrogante de qué ocurriría si alguien apagara la IA y nadie supiera cómo hacer las cosas.

\begin{note}
Reflexión: «Hay que aprender a hacerlo. Después usar las herramientas que te agilicen\ldots Es una flor de herramienta, pero hay que aprender a hacer las cosas, hay que aprender a darse cuenta cuando se está equivocando, cuando hay un error, cuando había una opción mejor.» Para pensar: ¿qué parte de mi trabajo podría revisar con criterio propio si la herramienta se equivocara?
\end{note}

% TODO: contenido incompleto o ambiguo en las notas originales (el tramo final de la grabación es poco inteligible: una estudiante relató cómo usa herramientas de IA para sus trabajos y cómo le respondieron; no es posible transcribirlo con fidelidad)

\section{Cuadro Cornell: IA y juicio profesional}

\begin{longtable}{@{}>{\raggedright\arraybackslash}p{0.27\textwidth}>{\raggedright\arraybackslash}p{0.69\textwidth}@{}}
\toprule
\textbf{Preguntas / palabras clave} & \textbf{Notas / respuestas} \\
\midrule
\endfirsthead
\toprule
\textbf{Preguntas / palabras clave} & \textbf{Notas / respuestas} \\
\midrule
\endhead
\cornellrow{¿Cómo usar la IA sin dejar de pensar?}{Como herramienta, sin perder la capacidad de cuestionar el resultado y de tomar la última decisión.}
\cornellrow{¿Por qué es necesario saber hacer el trabajo manualmente?}{Porque solo quien aprendió a hacerlo puede detectar los errores de la herramienta (por ejemplo, errores de sentido y párrafos omitidos en traducciones).}
\cornellrow{¿Qué riesgo plantea la delegación total en la IA?}{Que profesionales (médicos, arquitectos, abogados) hayan resuelto su formación con IA y no sepan hacer las cosas por sí mismos.}
\cornellrow{¿Qué aspecto señaló la docente como propio del derecho?}{El aspecto humano y sus matices, por ejemplo en cuestiones de familia o en el caso del 80\,\% de culpa concurrente, que la máquina no percibe.}
\midrule
\multicolumn{2}{@{}>{\raggedright\arraybackslash}p{0.96\textwidth}@{}}{\textbf{Resumen:} la IA es una herramienta valiosa, pero solo quien aprendió a hacer el trabajo por sí mismo puede detectar sus errores y conservar la decisión final. Esto vale para la traducción, el diseño y el derecho.} \\
\bottomrule
\end{longtable}

\end{document}
